%=============================================================================!
% 15.5  A SLOW DRIFT                                                          !
%=============================================================================!
\section{A slow drift}
\label{sec:cv-drift}

A run can settle and still not hold still. An integration step that is too
long lets the total energy creep, a thermostat set to the wrong target
moves the temperature slowly, and a transient that discarding the start
found in \cref{sec:cv-equilibration} did not remove leaves a slow slope. Each is a
trend under the fluctuations, and this section fits a straight line to a
run to find it, with an error that accounts for the run's memory.

\subsection*{A tyre that may be leaking}

A bicycle tyre is checked with a gauge every morning for a month. The
gauge scatters by \SI{0.05}{bar} from reading to reading, and a slow leak
of \SI{0.01}{bar} a day hides under that scatter from one day to the
next, but not over thirty days. The straight line that passes closest to
the thirty readings finds it, and the toolbox gives that line and the
error of its slope.

\begin{toolbox}[A straight line through scattered points]
Readings $A_i$ at the times $t_i$, $i = 1, \dots, n$, are to be described by
$A \approx a + \dot A\,(t - \bar t)$, with $\bar t$ the mean time. The line
of least squares chooses $a$ and the slope $\dot A$ to make the sum of
squared misses $S_{\mathrm m} = \sum_i\bigl[A_i - a - \dot A(t_i - \bar
t)\bigr]^2$ as small as it can be, where its rates of change with $a$ and
with $\dot A$ both vanish. With respect to $a$,
\begin{equation*}
   \frac{\partial S_{\mathrm m}}{\partial a} = -2\sum_i\bigl[A_i - a - \dot A(t_i -
   \bar t)\bigr] = -2\Bigl(n\bar A - na - \dot A\cdot0\Bigr) = 0 ,
\end{equation*}
since the deviations $t_i - \bar t$ add to zero, so $a = \bar A$. With
respect to $\dot A$, multiplying each miss by $t_i - \bar t$,
\begin{equation*}
   \frac{\partial S_{\mathrm m}}{\partial\dot A} = -2\sum_i(t_i - \bar t)\bigl[A_i -
   \bar A - \dot A(t_i - \bar t)\bigr] = 0 , \qquad
   \dot A = \frac{\sum_i(t_i - \bar t)(A_i - \bar A)}{S_{tt}} ,
\end{equation*}
with $S_{tt} = \sum_i(t_i - \bar t)^2$. The slope is a weighted sum of the
readings, $\dot A = \sum_iw_iA_i$ with $w_i = (t_i - \bar t)/S_{tt}$, since
$\sum_iw_i\bar A = 0$. If the misses about the true line are independent
with variance $\sigma_{\mathrm m}^2$, the variances of the terms add, and a
constant factor $w_i$ multiplies a variance by $w_i^2$
(\cref{sec:cv-mean}), so
\begin{equation*}
   \operatorname{var}(\dot A) = \sigma_{\mathrm m}^2\sum_iw_i^2
   = \frac{\sigma_{\mathrm m}^2}{S_{tt}^2}\sum_i(t_i - \bar t)^2
   = \frac{\sigma_{\mathrm m}^2}{S_{tt}} .
\end{equation*}
$\sigma_{\mathrm m}^2$ is estimated from the misses about the fitted line, their sum
of squares divided by $n - 2$, since two quantities have been fitted to
the same readings; for the thousands of samples of a run the difference
from $n$ is negligible.
\end{toolbox}

For thirty daily readings $S_{tt} = \SI{2247.5}{\day\squared}$, so the
slope is known to $0.05/\sqrt{2247.5} = 0.00105$\,bar a day. One
month of readings drawn by computer with the leak gives $\dot A =
(-0.0109 \pm 0.0012)$\,bar a day, nine errors from zero (the
misses of this month scatter by \SI{0.057}{bar}, hence 0.0012 rather than
0.00105): the leak is certain, though no single day shows it.

\subsection*{The slope of a correlated series}

The samples of a run are correlated, and so are their misses about the
line. The variance of the weighted sum is then, as in
\cref{sec:cv-correlation}, a double sum over pairs,
\begin{equation*}
   \operatorname{var}(\dot A) = \sigma_{\mathrm m}^2\sum_i\sum_jw_iw_j
   \operatorname{corr}(\lvert i - j\rvert)
   \approx \sigma_{\mathrm m}^2\sum_iw_i^2\Bigl[1 +
   2\sum_{k\ge1}\operatorname{corr}(k)\Bigr]
   = \frac{g\,\sigma_{\mathrm m}^2}{S_{tt}} ,
\end{equation*}
the approximation taking $w_j = w_i$ wherever $\operatorname{corr}$ is not
negligible, which holds when the run is much longer than its memory: across
a memory of about $g$ samples a weight changes by $g\dt/S_{tt}$, a fraction
$2g/n$ of the largest weight, $(t_{\mathrm f}/2)/S_{tt}$. The slope's
error carries the same factor $\sqrt g$ as the mean's, with $g$ found from
the misses. For $n$ samples evenly spaced by $\dt$, $S_{tt} = n(n^2 -
1)\dt^2/12$ (\cref{ex:cv-stt}), close to $nt_{\mathrm f}^2/12$ for a run
of length $t_{\mathrm f} = n\dt$, so with $\tau_{\mathrm{int}} = g\dt/2$,
\begin{equation*}
   \operatorname{var}(\dot A) \approx
   \frac{12\,g\,\sigma_{\mathrm m}^2\dt}{t_{\mathrm f}^3}
   = \frac{24\,\sigma_{\mathrm m}^2\tau_{\mathrm{int}}}{t_{\mathrm f}^3} :
\end{equation*}
the error of a slope falls as $t_{\mathrm f}^{-3/2}$, and a run four times
as long can detect a slope eight times smaller. The ratio of the slope to
its error then tests for a drift: for a series with none it is Gaussian
with unit spread, and lies beyond $\pm2$ in 4.6\% of runs
(\cref{sec:sm-probability}). Applied to 1000 series of 5000 samples of the
kind in \cref{fig:cv-ou}, with $g = 39$ and no drift, the
test \texttt{drift\_test} of \texttt{mdlab.analysis.stats} finds a drift
in 0.051 of them. A slope beyond two errors is evidence of drift worth investigating; within them, the run shows none that it can detect.

\subsection*{The time step}

Chapter~9 judged a time step by the spread of the total energy at fixed
energy, which should fall as $\dt^2$ and stay small beside the spread of
the kinetic energy, and by the absence of a drift (\cref{sec:in-stability}).
The drift test makes the last of these quantitative.
\Cref{fig:cv-drift} follows the liquid at fixed energy for
\SI{1}{\nano\second} at each of six steps. The spread of the total energy
over that of the kinetic energy is 3.93 times larger at 10 than at
\SI{5}{\femto\second} and 4.12 times larger at 20 than at
\SI{10}{\femto\second}, the factor 4 of $\dt^2$, but 3.93 times larger at
30 than at \SI{20}{\femto\second}, where $\dt^2$ gives 2.25, and 6.18 times
larger at 35 than at \SI{30}{\femto\second}, where it gives 1.36: beyond
\SI{20}{\femto\second} the step is leaving the range in which its error is
a small, smooth correction. The drift per atom per nanosecond is $0.00
\pm 0.01$ and $-0.04 \pm 0.05$\,\si{\micro\electronvolt} at 5 and
\SI{10}{\femto\second}; $-0.66 \pm 0.17$\,\si{\micro\electronvolt} at
\SI{20}{\femto\second}; $-54 \pm 29$\,\si{\micro\electronvolt} at
\SI{30}{\femto\second}; $+642 \pm 69$\,\si{\micro\electronvolt} at
\SI{35}{\femto\second}; and at \SI{40}{\femto\second} the energy climbs by
\SI{11}{\milli\electronvolt} per atom per nanosecond over the
\SI{363}{\pico\second} it takes to gain \SI{1}{\electronvolt} in all, and
the run blows up at \SI{557}{\pico\second}.

Whether a drift matters depends on what the run measures. At fixed energy
a drift $\dot E$ moves the temperature at the rate $\dot E/C_V$, and over
a run of length $t_{\mathrm f}$ moves its mean by half of $\dot
Et_{\mathrm f}/C_V$. Asking that this be smaller than the mean's own error
$\sigma_{\bar T}$,
\begin{equation*}
   \frac{\lvert\dot E\rvert}{N} < \frac{2\sigma_{\bar T}}{t_{\mathrm f}}\,
   \frac{C_V}{N} .
\end{equation*}
For this liquid $C_V/N = 2.347\kB$, the mean of the five single-run
estimates of \cref{sec:sm-kinetic} (the slope of $E$ against $T$ gives
2.376, within the 3\% errors of those runs, \cref{sec:cv-observables}), and $\sigma_{\bar T} =
\SI{0.0513}{\kelvin}$ over \SI{1}{\nano\second} at fixed
energy, so the limit for a run of \SI{1}{\nano\second} is
\SI{20.8}{\micro\electronvolt} per atom per nanosecond (dotted in
\cref{fig:cv-drift}(b)). The drift at \SI{20}{\femto\second} is significant,
3.8 errors from zero, and harmless: it moves the temperature by
\SI{0.0033}{\kelvin} in a nanosecond. The drift at \SI{30}{\femto\second}
lies 1.8 errors from zero and 1.1 above the limit, so the run can show it
neither significant nor harmful; at \SI{35}{\femto\second} it moves the
temperature by \SI{3.2}{\kelvin} per nanosecond. Significance grows with the length of a run, which can detect
ever smaller drifts; importance is set by the quantity measured. For this liquid \SI{20}{\femto\second} passes Chapter~9's tests
of the spread, and its drift, though significant, is harmless for a run of
\SI{1}{\nano\second} and at the limit for one of \SI{10}{\nano\second}
(\cref{ex:cv-drift-limit}); \SI{10}{\femto\second}, the step used
throughout, passes them with a wide margin and shows no drift the test can
detect.

\begin{figure}[htb]
\centering
\includegraphics{ch15_convergence/drift}
\caption{A slow drift: the liquid at fixed energy for \SI{1}{\nano\second}
at each step. (a) The total energy per atom less its start, at 20 (teal),
30 (ochre), 35 (oxblood) and \SI{40}{\femto\second} (grey, until it passes
\SI{1}{\electronvolt} in all). (b) The size of the drift per atom per
nanosecond, with its error, against the step, and the largest drift that
moves the mean temperature of a \SI{1}{\nano\second} run by less than its
error (dotted). (c) The spread of the total energy over that of the
kinetic energy, against $\dt^2$ through the point at \SI{5}{\femto\second}
(dashed).}
\label{fig:cv-drift}
\end{figure}

\begin{keyidea}
A drift is the slope of the line of least squares through a run, with the
error $\sqrt{g\sigma_{\mathrm m}^2/S_{tt}}$ that the run's memory requires; a slope
beyond two errors is evidence of drift under this test. Whether it matters depends on what is
measured: at fixed energy it is harmless while it moves the mean
temperature by less than its error. A step is acceptable while the
energy's spread grows as $\dt^2$ and its drift is harmless.
\end{keyidea}

\notebook{15}{hide a leak under noise and test for it}

\begin{selfcheck}
A run of \SI{500}{\pico\second} gives a drift of $(3 \pm 2)$ in some unit.
Is it a drift, and what would a run of \SI{2}{\nano\second} with the same
true slope be expected to show?
\end{selfcheck}
